% =====================================================================================
% theory_si.tex -- SI appendix: assumptions, derivations and mapping for the main-text theory section
% (paper/theory_section.tex). Written 2026-09-25 AFTER the results were seen, and revised the same day after an
% adversarial review (THEORY_NOTES.md section 10 lists every change and its reason).
% New references: paper/refs_theory_pending.bib (Crossref-verified); all other keys are in refs_merged.bib.
% Uses \ci (defined in si.tex and main.tex) and longtable (loaded in si.tex).
% Every number is copied from a RESULT write-up or interim table named in the text and was re-checked against it on
% 2026-09-25; the [P] numbers of the Proposition 2 subsection were recomputed from data/interim/field_vs_generic.csv
% (scripts/28 output). The SpringRank counterexample and the gender-missingness shares were recomputed on
% 2026-09-25 (THEORY_NOTES.md section 10). The binding texts of the registered tests (B0 and amendment A1) are in
% THEORY_NOTES.md section 9 and in the docstring of scripts/75_theory_tests.py.
% \theoryprereg: where the binding texts are deposited. The author must replace the default with the public commit
% hash or DOI before any test is run.
\providecommand{\theoryprereg}{the docstring of \texttt{scripts/75\_theory\_tests.py} (public deposit pending;
the commit hash or DOI will be given here before any test is run)}
% =====================================================================================
\section{A model of two audiences: assumptions, derivations and mapping}\label{si:theory}

\subsection*{Status of this appendix}

This appendix states the assumptions of the model in the main-text theory section, proves its propositions, and
maps every empirical result of the paper to a proposition. The model was written on 25 September 2026, after all
results in the main text and the SI had been seen. It synthesises three candidate formulations: an
industrial-organisation model of reputation, an employer-learning model and a measurement model. Claims that did
not survive two rounds of adversarial review were dropped or restated. The mapping in
Section~\ref{si:theory:map} was organised after the results were seen. It is a check of post hoc consistency and
is not evidence for the model. The model's out-of-sample content is the registered tests of
Section~\ref{si:theory:tests}, none of which has been run.

The two construct readings are argued as interpretive arguments \citep{kane2013validating}. Each has a
definition (the audience, its product and its revealing choice), a warrant (Propositions 1a and 1b), and
rebuttal conditions (the failure conditions listed with each assumption). The prospective backing is the
registered tests; the mapping below is post hoc consistency, not backing. The readings also meet the criterion of
\citet{borsboom2004concept}: in the model, variation in the audience's valuation produces variation in the
indicator, through hiring flows or through wage offers. The tests extend the nomological network of
\citet{cronbach1955construct}; T2 has the convergent--discriminant form of \citet{campbell1959convergent}, but one
half of its contrast is already known.

Each derivation carries one of three tags:
\begin{itemize}
\item \textbf{[S]} (structural): holds for every admissible parameter value under the stated assumptions.
\item \textbf{[B]} (baseline): holds under an added assumption that is named, which a named alternative contradicts.
\item \textbf{[P]} (post hoc): a calibration computed from results already in hand. It could have failed, but it
was computed after the data were seen.
\end{itemize}
These tags describe derivations. How results fit the model is classified separately in
Section~\ref{si:theory:map}.

% -------------------------------------------------------------------------------------
\subsection{Setup}\label{si:theory:setup}

\paragraph{Units and products.} In field $f$, $\mathcal I_f$ is the set of institutions that grant the
bachelor's degree in $f$ and have a ranked doctoral department in $f$. A program is $p=(i,f)$. Field $f$ belongs
to an organisational school $d(f)$ (for example engineering, business, or arts and sciences). A department sells
two products under the name $i$: doctoral graduates, bought by the field's hiring committees (the
\emph{academy}), and bachelor's graduates, bought by employers (the \emph{market}).

\paragraph{Latent structure.} The academy's valuation of department $(i,f)$ and the mean productivity of its
bachelor's graduates are
\begin{equation}\label{eq:T-latent}
s_{if}=\alpha_fa_i+r^S_{i\,d(f)}+r_{if},\qquad
\theta_{if}=\mu_f+\beta_fx_i+y^S_{i\,d(f)}+y_{if},\qquad \eta_k=\theta_{if}+\varepsilon_k .
\end{equation}
Here $\eta_k$ is the productivity of graduate $k$, $\varepsilon_k\sim N(0,\sigma^2)$, and $\theta_{if}$ bundles
selection, peers, value added and networks. $\alpha_f$ is the weight of the institution component in field $f$'s
academic valuation, and $\beta_f$ is how strongly field $f$'s labour market rewards the institution component of
graduates' productivity. Within field $f$, $r^S+r$ is defined as the residual of $s$ on $a$, and $y^S+y$ as the
residual of $\theta$ on $x$, over $\mathcal I_f$; so $\mathrm{Cov}_f(a,r^S+r)=\mathrm{Cov}_f(x,y^S+y)=0$ by
construction. The two sides are linked within level:
\[
\mathrm{Corr}(a,x)=\chi,\qquad \mathrm{Corr}(r^S,y^S)=\kappa^S_d,\qquad \mathrm{Corr}(r_{\cdot f},y_{\cdot f})=\kappa_f .
\]
The variances are $\sigma_a^2,\sigma_{rS}^2,\sigma_r^2$ on the academic side and $\sigma_x^2,\sigma_{yS}^2,\sigma_y^2$
on the graduate side.

\paragraph{A-orth (no covariance across levels).} $\mathrm{Cov}_f(r^S+r,\,x)=0$,
$\mathrm{Cov}_f(a,\,y^S+y)=0$, and the school and department deviations are uncorrelated with institution-level
growth $g_i(t)$. The projections above do not deliver these; they are assumptions. A-orth fails if departments
are relatively strong at highly selective institutions ($\mathrm{Cov}(r,x)>0$), if programs at high-status
institutions draw relatively able students into the major ($\mathrm{Cov}(a,y)>0$), or if strong departments sit at
institutions whose graduates' pay grows faster.

\paragraph{Indicators.} The indicators are
\begin{itemize}
\item $F_{if}=s_{if}+e^F_{if}$, the negated published field rank (SpringRank), whose reliability
$\mathrm{rel}_F$ is bracketed at 0.81--0.96 (\texttt{scripts/56});
\item $G_i=a_i+\bar r_i+e^G_i$, the negated published academia-wide rank, with $\mathrm{rel}_G\approx0.99$;
\item $Y_{if}(t)=m_{if}(t)+e^Y_{if}(t)$, the program's median log earnings, with reliability 0.78 (median over
fields) for Scorecard and 0.74--0.76 for PSEO (\texttt{scripts/73}).
\end{itemize}
In $G_i$, $\bar r_i$ is the hire-weighted mean of the institution's school and department deviations. Indicator
errors are classical: mean zero, mutually independent, and independent of the latent components.

\paragraph{A0 (Gaussian copula).} Within a field, the normal scores of the latent components and the indicators
are jointly Gaussian. Spearman's $\rho$ and the Pearson $r$ of the normal scores are then linked by
$\rho=\tfrac{6}{\pi}\arcsin(r/2)$ \citep{kruskal1958ordinal}, a strictly increasing map. Statements about signs
and orderings hold on either scale. Magnitudes (the thresholds, decompositions and $\cos_\Sigma$ values below) are
statements about normal scores, carried to $\rho$ by this map. A rank-based partial Spearman correlation need not
have the sign of the latent partial correlation when that is near zero. Where the latent sign matters (test T2),
partial correlations are also computed from the arcsine-transformed pairwise $\rho$.

% -------------------------------------------------------------------------------------
\subsection{The academy's valuation (Proposition 1a)}\label{si:theory:academy}

Let $A_{uv}$ be the number of field-$f$ faculty employed at $v$ whose PhD is from $u$, with self-hires ($u=v$)
excluded. Let $d^{\rm out}_u=\sum_vA_{uv}$ and $d^{\rm in}_u=\sum_vA_{vu}$.

\paragraph{Assumptions.}
\begin{itemize}
\item \textbf{H0 (volume neutrality).} Write $\mathbb E A_{uv}=p_uh_v\,\varphi(s_u-s_v)\,\zeta_{uv}$, with $p_u$
the production of PhDs at $u$, $h_v$ the hiring at $v$, $\varphi$ increasing, and $\zeta_{uv}=\zeta_{vu}$
collecting symmetric factors (geography, subfield overlap). H0: the production-to-hiring ratio $p_u/h_u$ is the
same for all departments. Then $\mathbb E A_{uv}/\mathbb E A_{vu}=\varphi(\Delta)/\varphi(-\Delta)$ with
$\Delta=s_u-s_v$, so the direction of exchange depends on the valuation alone. Where $p_u/h_u$ varies, the ratio
is multiplied by $(p_u/h_u)/(p_v/h_v)$, and H0 fails to the extent that this factor correlates with $s$ or reverses
the direction.
\item \textbf{H1 (vertical sorting).} For every pair with $\mathbb E(A_{uv}+A_{vu})>0$,
$s_u>s_v\Rightarrow\mathbb E A_{uv}>\mathbb E A_{vu}$.
\item \textbf{H2 (many hires).} $\mathbb E A_{uv}=n\,\alpha_{uv}$ with $\alpha$ fixed and $n\to\infty$; the set
of institutions is fixed.
\item \textbf{H3 (other drivers do not set direction).} Geography, dual careers, subfield fit and advisor networks
can change how many hires a pair exchanges, but they enter $\mathbb E A_{uv}$ and $\mathbb E A_{vu}$ alike
(through $\zeta_{uv}$) or independently of $s$, so they do not reverse H1.
\item \textbf{H4 (stationarity).} $s$ is constant over the decades in which the observed faculty were hired.
\end{itemize}
H0 implies H1, but H1 can hold without any valuation. If $\varphi$ is constant
(no valuation, random matching), $\mathbb E A_{uv}\propto p_uh_v$, and every pair satisfies H1 with departments
ordered by $p_u/h_u$. H0 is what makes the direction of exchange informative about $s$. The public edge list
reports $A_{uv}$ only when $A_{uv}\ge2$, and the published ranks are estimated on the census.

\paragraph{Consistency of the minimum-violation ordering [S].} For an ordering $\pi$, the expected number of
hires that move up $\pi$ is $V(\pi)=\sum_{(u,v):\,v\succ_\pi u}\mathbb E A_{uv}$. $V$ is a sum over unordered
pairs. For a pair with $s_u>s_v$, an ordering that places $v$ above $u$ adds $\mathbb E A_{uv}$ to $V$, while
the $s$-ordering $\pi_s$ adds $\mathbb E A_{vu}$. Hence, by H1,
\[
V(\pi)-V(\pi_s)=\sum_{\text{pairs inverted by }\pi}(\mathbb E A_{uv}-\mathbb E A_{vu})\ge0 ,
\]
with equality only if $\pi$ inverts no pair that exchanges faculty. So $\pi_s$ minimises $V$, and every minimiser
agrees with $s$ on every exchanging pair. A minimiser is a total order, so it also agrees with $s$ on every pair
linked by a chain of exchanging pairs that is monotone in $s$; pairs linked by no such chain are left unordered
by the data, and most institution pairs in a field never exchange. By H2, $V_n(\pi)/n\to V(\pi)$ for each of the
finitely many orderings. The gap $V(\pi)-V(\pi_s)$ is bounded away from zero for every $\pi$ that inverts an
exchanging pair, so the sample minimiser \citep{clauset2015systematic} agrees with $s$ on those pairs with
probability tending to one. $\square$

This is a consistency property, close to a restatement of H1. The substantive content is H0 and H1, which define
the construct and can fail.

\paragraph{What $F$ is: SpringRank, not the minimum-violation ordering.} The published $F$ is the SpringRank
ground state, the minimiser of $H(s)=\tfrac12\sum_{uv}A_{uv}(s_u-s_v-1)^2$ (with a small ridge). Its link to $s$
rests on its own generative model, which is stronger than H1.
\begin{itemize}
\item \emph{Generative model [S].} If $A_{uv}\sim\mathrm{Poisson}(c\exp\{-\tfrac{\beta}{2}(s_u-s_v-1)^2\})$
independently, with one constant $c$, then $\mathbb E A_{uv}/\mathbb E A_{vu}=e^{2\beta\Delta}$ and, given the
pair total $T_{uv}$, $A_{uv}\mid T_{uv}\sim\mathrm{Bin}\{T_{uv},\Lambda(2\beta\Delta)\}$ with
$\Lambda(z)=1/(1+e^{-z})$. As \citet{debacco2018physical} show, this is the Bradley--Terry--Luce comparison
\citep{bradley1952rank} with strengths $e^{2\beta s_u}$, and the maximum-likelihood ranks of this model approach
the SpringRank ground state only in the limit of strong hierarchy (large $\beta$). The model has no degree
heterogeneity.
\item \emph{Disagreement under H1 alone.} Take three departments with $A_{uv}=1$, $A_{vu}=0$, $A_{vw}=100$,
$A_{wv}=0$, $A_{uw}=100$ and $A_{wu}=99$. H1 holds for $u\succ v\succ w$, and the minimum-violation ordering is
$u\succ v\succ w$ (99 upward hires; every other ordering has at least 100). The SpringRank ground state
(\texttt{src/ar\_pipeline/springrank}, $\alpha=10^{-3}$) gives $s_v-s_w=0.98$ and $s_u-s_w=0.015$, so $v\succ
u\succ w$. We recomputed this on 25 September 2026.
\item \emph{Net exporters rank higher [S].} With $\alpha\to0$, the first-order condition of $H$ is
\[
s_u=\bar s_u+\frac{d^{\rm out}_u-d^{\rm in}_u}{d^{\rm out}_u+d^{\rm in}_u},\qquad
\bar s_u=\frac{\sum_v(A_{uv}+A_{vu})\,s_v}{d^{\rm out}_u+d^{\rm in}_u}.
\]
A department sits above the exchange-weighted mean of its partners by its net-export ratio. Where H0 fails and
$p_u/h_u$ varies, net exports reflect $p_u/h_u$ as well as $s_u$, so at equal valuation SpringRank ranks the
department that produces more PhDs per faculty slot higher. In the no-valuation null above, SpringRank
ordered departments exactly by $p_u/h_u$ in a numerical check on a random instance. Faculty production is highly
unequal \citep{wapman2022quantifying}.
\end{itemize}
$F$ is therefore a department's net-placement position. It reads as the academy's valuation only under H0 (test
T0 audits this; T0 also reports the Spearman correlation of $F$ with a minimum-violation ordering computed on the
public edges).

\paragraph{Remarks.}
\begin{enumerate}
\item \emph{What H1 is.} H1 is not a lemma. It is the revealed-preference content of the construct: a
department's standing is revealed by the net direction in which hires move between it and each department it
trades with, net of volume (H0). The comparison with \citet{avery2013revealed}, who read colleges' standing from
the choices of students admitted to both, is an analogy, not an equivalence. There the same chooser picks between
two options. Here $A_{uv}$ and $A_{vu}$ are different events, involving different candidates and different
hiring slots, faculty hiring is two-sided matching shaped by candidates' location and dual-career constraints,
and PhDs who were not placed are not observed, so the rejected alternatives are missing. A sufficient
micro-foundation, which we state but do not prove, has three parts: committees rank candidates by the valuation of
their PhD department plus idiosyncratic noise; candidates rank employers the same way; and production and hiring
volumes are proportional (H0). With common rankings on both sides the stable matching \citep{gale1962college} is
positively assortative, and net flows run down the valuation.
\item \emph{Two roles, one index.} Under a single index, the placement position of $u$ (the mean standing of the
employers of $u$'s PhDs) and its attraction position (the mean standing of the origins of $u$'s faculty) both
increase with $s_u$, so they order departments alike up to noise and boundary effects at the top and bottom of
the hierarchy. If committees valued a department as a producer on one index and as an employer on another, the two
positions would diverge. Test T1 checks this against the value that SpringRank's own generative model produces,
because boundary effects make the two positions disagree even under one index. A world in which $F$ is only a
position in a hierarchical network, or only placement power, predicts the same agreement, so T1 does not separate
valuation from network position.
\item \emph{Status or belief.} The revealed ordering is status in a deference network
\citep{burris2004academic,podolny1993status,sauder2012status}. PhD origin shapes placement beyond productivity
\citep{long1979entrance}, so status may be partly self-reinforcing. $F$ reads as a belief about the quality of a
department's PhDs only if one index drives both roles, sorting is vertical and volume does not set direction. If
T1 fails, only that reading is withdrawn: $F$ is then described as the ordering revealed by net placement, its
position in the hiring network. How a producer is valued depends on the categories its audience uses
\citep{zuckerman1999categorical}, which is one reason the two audiences are modelled separately.
\item \emph{Academia-wide rank.} To first order, SpringRank on the network pooled over fields gives
$G_i\approx a_i+\bar r_i$ with $\mathrm{Var}(\bar r)=O(\bar\sigma_r^2/K_{\rm eff})$, so $G$ is the academy's
valuation of the name, subject to the same H0. That $G$ is the institution's ``brand'' is an interpretation;
nothing requires employers to see $G$.
\item \emph{Failure conditions.}
\begin{itemize}
\item H0 fails where production-to-hiring ratios differ across departments in a way that sets the direction of
exchange (test T0).
\item H1 fails if matching is horizontal in a way correlated with $s$, for example by subfield
\citep{laberge2022subfield}, or if the academy values departments on more than one index (T1).
\item H2 fails in thin fields. Because single-hire pairs are suppressed, $\mathrm{rel}_F$ can only be bracketed.
\item H3 fails if steering by subfield, geography or dual careers differs by candidates' gender (T1(b) audits
this).
\item H4 fails for departments that rose or fell quickly. The faculty stock behind the 2011--2020 ranks was hired
over decades, and our graduates are the cohorts of 2001--2019.
\item The Wapman-to-CIP crosswalk can be wrong.
\end{itemize}
\end{enumerate}

% -------------------------------------------------------------------------------------
\subsection{The market's valuation (Proposition 1b)}\label{si:theory:market}

\paragraph{Assumptions.}
\begin{itemize}
\item \textbf{M1 (competitive pricing on a record).} $w_k(t)=\hat{\mathbb E}[\eta_k\mid\Omega_k(t),H_t]$:
employers pay their posterior expectation of the graduate's productivity given the credential and the graduate's
own signals $\Omega_k(t)$, and a record $H_t$ of the performance of past graduates that employers have attributed
to units, under a normal prior over unit means \citep{farber1996learning}. Attribution is costly (bounded
attention), or labour markets are segmented, so the count $n_u$ that enters a unit's record is the number of its
graduates that an employer, or a local labour market, has observed and attributed to it, not the number ever
graduated.
\item \textbf{M2 (credential observed).} $\Omega_k(0)$ contains the institution and the major.
\item \textbf{M3 (no schedule).} Pay is not set by a scale that is independent of the posterior expectation.
\item \textbf{M4 (coverage and shape).} The program median covers graduates who work in the employer market, and
the within-program distribution of log earnings is symmetric, so the median equals the mean of log pay. For a
mixture of scheduled and market pay the median is not linear in the scheduled share; Proposition 5 uses the
mean of log pay as an approximation.
\end{itemize}
Under M1 with a public, costless record, $n_u$ would be every graduate ever observed, and with hundreds of
graduates per program $\omega_f$ would be close to one unless $\sigma_y^2$ is tiny. A small $\omega_f$ therefore
needs costly attribution or segmented records, and a tiny $\sigma_y^2$ is itself a content explanation (little
departmental variation in graduates' productivity to price).

\paragraph{Information.} At hiring, employers' belief about a graduate of $p$ is
\[
R^E_p=\mu_f+\beta_f\hat x_i+\omega^S_d\,y^S_{i\,d}+\omega_f\,y_{if},\qquad \hat x_i=\omega_ix_i+(1-\omega_i)\,\mathbb E[x\mid P_i],
\]
where $P_i$ is a public, non-academic signal attached to the name (name recognition, say). Each weight has the form
$\omega_u=n_u\tau_u^2/(n_u\tau_u^2+\sigma^2)$ (Section~\ref{si:theory:attrib}). Signals about the individual
arrive at hiring (precision $h_b$) and each year (precision $h_\varepsilon$); with prior precision $h_0$, the
weight on the graduate's own record is
\[
\lambda_t=b+(1-b)\frac{t}{t+k},\qquad b=\frac{h_b}{h_0+h_b},\quad k=\frac{h_0+h_b}{h_\varepsilon} .
\]

\paragraph{Derivation of the program wage [S].} Each wage is
$w_k(t)=(1-\lambda_t)R^E_p+\lambda_t(\eta_k+\bar\nu_{kt})$, where $\bar\nu_{kt}$ is the mean of $k$'s signal
noise. Averaging over the program's graduates removes $\varepsilon_k$ and $\bar\nu_{kt}$. Adding status-dependent
growth $g_p(t)$ gives
\begin{equation}\label{eq:T-wage}
m_p(t)=\mu_f(t)+\beta_f\big[(1-\lambda_t)\hat x_i+\lambda_tx_i\big]+\Lambda^S_{dt}\,y^S_{i\,d}+\Lambda_{ft}\,y_{if}+g_p(t),
\qquad \Lambda_{ut}=\omega_u+(1-\omega_u)\lambda_t .
\end{equation}
With $\hat x_i=x_i$ (Lemma~B) this is the program wage of Proposition~1(b) in the main text, plus the school term.
$\square$

\paragraph{Lemma A (invariance) [S].} If employers' posterior means of the unit means equal the unit means
($\omega_i=\omega^S_d=\omega_f=1$), then $m_p(t)-g_p(t)-\mu_f(t)=\theta_p-\mu_f$ for every $t$.

\emph{Proof.} Then $R^E_p=\theta_p$, and $(1-\lambda_t)\theta_p+\lambda_t\theta_p=\theta_p$. $\square$

This is the result of \citet{farber1996learning}, that the return to a variable employers observe at hiring does
not change with experience, applied to program means. Learning about individuals moves each graduate's wage but
not the program mean, unless employers are uncertain about the program's own mean. The lemma is conditional on
the posterior mean equalling $\theta_p$; it does not follow from M1 and M2 alone.

\paragraph{Corollaries [S].}
\begin{itemize}
\item \emph{Entry.} With $b=0$, $m_p(0)=R^E_p$: the program's entry wage is employers' belief about the program,
its employer reputation in the sense of \citet{macleod2015reputation}, selection included.
\item \emph{Name.} By Lemma~B, $\omega_i\to1$ for institutions whose graduates employers see often. At the name,
employer reputation then equals the realised mean productivity of the institution's graduates, and the two differ
only below the name, by $(1-\omega)(1-\lambda_t)$ times the school or department deviation.
\item \emph{Employer reputation in the learning sense.} The component that \citet{macleod2017big} identify, a
school-mean premium with graduates' own ability held fixed that learning erodes, requires graduates' own ability
to be observed. The US data do not observe it. The UK data observe prior-attainment bands, and within a band
selectivity keeps $+0.252$ beyond $G$ (\texttt{scripts/62}); test T4 registers whether that premium decays.
\end{itemize}

\paragraph{Failure conditions.}
\begin{itemize}
\item M1 fails under monopsony \citep{staiger2010there}, rent sharing and compensating differentials, for example
when graduates of esteemed programs take lower-paid academic or public jobs.
\item M3 fails where pay follows schedules (Proposition~5).
\item M4 fails in several ways. Graduate study depresses y1 medians. Scorecard covers Title IV recipients only.
PSEO coverage falls with prestige at y1. Hours, labour supply (part-time work, career breaks) and the gender
composition of a program move its median; in the UK, standardising for sex as well as prior attainment lowers
the share of coupling that survives to 66\% (\texttt{scripts/72}). Medians hide the tails.
\item \emph{Location.} Raw pay contains local price levels, which lie outside the construct, and access to
high-wage labour markets, which is a quantity. The only location controls so far are at the institution's own
state: the state earnings level (\texttt{scripts/55}) and the field-specific earnings component of the
institution's state, which moves mean coupling from $+0.419$ to $+0.415$ (\texttt{scripts/69}). No result nets out
where graduates work. Test T7 registers that control; until it is run, construct claims about pay are not netted
of destination location.
\item \emph{Sector composition.} If placement into better-paid sectors is inside the pay construct, the construct
is the market's valuation including allocation. Reading between-sector coupling as employers' valuation requires
sector access to be rationed by employers; supply-side sorting (graduates' preferences, family background, peer
norms) breaks this. Test T8 registers the decomposition.
\item \emph{Selection.} At the name the construct includes selection. Much of the earnings premium of selective
colleges reflects whom they admit \citep{dale2002estimating,dale2014estimating}, and their causal effects
concentrate in top-tail outcomes (top jobs and incomes) that medians miss \citep{zimmerman2019elite,chetty2023diversifying}.
\item \emph{The umbrella.} If the program is printed on the CV and performance records are public and costless,
attribution is free ($s_f\approx1$) and $\omega_f\approx1$ (see the remark after M4).
\end{itemize}

% -------------------------------------------------------------------------------------
\subsection{Information and attribution}\label{si:theory:attrib}

\paragraph{Lemma B (umbrella precision) [S].} An employer has observed $n^e_{ig}$ graduates of program $(i,g)$.
The name's record pools all of them, $n^e_i=\sum_gn^e_{ig}$, and its mean is $x_i$ plus
$u_i=\sum_g(n^e_{ig}/n^e_i)\,(y^S_{i\,d(g)}+y_{ig})$, a term common to all of the institution's programs. Each
observed graduate is attributed to the department with probability $s_g$, and department deviations are learned
as deviations from the name's mean, from $n^e_{ig}s_g$ observations. Hence
\[
\omega_i=\frac{n^e_i\tau_i^2}{n^e_i\tau_i^2+\sigma^2},\qquad
\omega_f=\frac{n^e_{if}s_f\tau_y^2}{n^e_{if}s_f\tau_y^2+\sigma^2}.
\]
Since $n^e_i\ge n^e_{if}s_f$, $\omega_i\ge\omega_f$ whenever $\tau_i^2\ge\tau_y^2$; $\omega_i\to1$ as
$n^e_i\to\infty$; and $\omega_f\to0$ as $s_f\to0$. Deviations that are never attributed enter pay only through
$u_i$, which institution fixed effects absorb. $\square$

The counts are per employer or per local market (M1). With $\tau^2/\sigma^2=0.04$, a record of 1,000 graduates
gives $\omega=0.976$, so an institution whose graduates a given employer sees rarely is the only place where
$\omega_i$ can be well below one.

\paragraph{Attribution choice.} Tracking department histories in field $f$ lowers an employer's prediction error
by $\omega(n_{if})\sigma_y^2$ per hire. The employer tracks them if and only if
$S_f\,\omega(n_{if})\,\sigma_y^2\ge c_f$, where $S_f$ is the stake (hires in the field times the sensitivity of pay
to prediction error) and $c_f$ the cost of attribution. Public program-level signals lower $c_f$: program
rankings, admission controlled by the department, accreditation, a separately branded school, and published
program-level earnings. The same condition with school counts governs $s^S_d$. The pooling is analogous to
umbrella branding \citep{wernerfelt1988umbrella}. The mechanism here is buyers' inference across the products of
one brand \citep{choi1998brand,cabral2000stretching}, not a firm's decision to signal through the brand, and
reputations shared by the members of a group are collective reputations in the sense of \citet{tirole1996theory}.

% -------------------------------------------------------------------------------------
\subsection{Proposition 2: decomposition, $F$ against $G$, and what coupling measures}\label{si:theory:decomp}

\paragraph{Decomposition [S].} Take growth to vary across institutions only, $g_p(t)=g_i(t)$, as in the main
text. Substituting (\ref{eq:T-latent}) and (\ref{eq:T-wage}) with $\hat x_i=x_i$, and using classical errors,
\begin{align}
\mathrm{Cov}\{F_{if},m_{if}(t)\}={}&\alpha_f\beta_f\chi\sigma_a\sigma_x+\Lambda^S_{dt}\kappa^S_d\sigma_{rS}\sigma_{yS}
+\Lambda_{ft}\kappa_f\sigma_r\sigma_y+\alpha_f\mathrm{Cov}\{a_i,g_{i}(t)\}\label{eq:T-decomp}\\
&+\underbrace{\beta_f\mathrm{Cov}(r^S+r,x)+\alpha_f\mathrm{Cov}(a,\Lambda^S y^S+\Lambda y)+\mathrm{Cov}(r^S+r,g_t)}_{=0\text{ under A-orth}} .\nonumber
\end{align}
\begin{enumerate}
\item[(i)] $\mathrm{Cov}(G,m_t)=\beta_f\chi\sigma_a\sigma_x+\mathrm{Cov}(a,g_t)+O(\mathrm{Var}\,\bar r)$ under A-orth.
\item[(ii)] Under A-orth, the partial covariance of $F$ and $m$ given $G$ is the school term plus the department
term; given $G$ and the school's standing, it is the department term alone. Without A-orth it also contains the
cross-level terms, so the cross-institution partial of $F$ given $G$ identifies $\Lambda_{ft}\kappa_f$ only under
A-orth. Cross-level terms of this kind would make the cross-institution partial exceed what the
within-institution design recovers.
\item[(iii)] \emph{Within institutions.} Institution fixed effects absorb what is common to an institution's
programs: $a_i$, the field-average reward $\bar\beta x_i$, institution-level growth and $u_i$. They do not absorb
field-specific rewards $(\beta_f-\bar\beta)x_i$, so identification needs field-specific slopes on institution
traits, as in the main specification of \texttt{scripts/58} (field $\times$ \{SAT, admission rate, Pell share,
$G$\}). Those slopes remove $(\beta_f-\bar\beta)x_i$, and the $\mathrm{Cov}(r,x)$ part of A-orth, only as far as
SAT, admission rate, Pell share and $G$ span $x$. Given that, the slope of $m$ on $F$ is
\[
\mathrm{rel}^w_F\,\frac{\Lambda^S\kappa^S\sigma_{rS}\sigma_{yS}+\Lambda_{f}\kappa_f\sigma_r\sigma_y}{\sigma_{rS}^2+\sigma_r^2},
\qquad \mathrm{rel}^w_F=1-\frac{1-\mathrm{rel}_F}{1-R^2_{F\mid W}} .
\]
The school term enters only for institutions with programs in several schools. With $R^2_{F\mid W}\approx0.72$
(\texttt{scripts/58}), $\mathrm{rel}^w_F$ is 0.32--0.86 across the reliability bracket. The corrected slopes of
\texttt{scripts/58} ($+0.0297$ and $+0.0098$ against $+0.0084$) imply 0.28 and 0.86. This agreement is an
identity, not a test.
\end{enumerate}

\paragraph{The $F$-against-$G$ identity [S].} For Spearman correlations computed on the same institutions,
Pearson algebra on ranks gives exactly
\[
c_F=r_{FG}\,c_G+\sqrt{1-r_{FG}^2}\;\mathrm{sr}_F,\qquad
\mathrm{sr}_F=\frac{c_F-r_{FG}c_G}{\sqrt{1-r_{FG}^2}}=\mathrm{Corr}(Y,F^{\perp G}) .
\]
Hence $c_F>c_G\iff\mathrm{sr}_F>c_G(1-r_{FG})/\sqrt{1-r_{FG}^2}=c_G\tan(\vartheta/2)$ with
$\cos\vartheta=r_{FG}$. The identity holds for any three variables; the model enters only in reading
$\mathrm{sr}_F$. The part of $F$ that $G$ lacks is signal for the academy and noise for the market unless
$\kappa_f\Lambda_{ft}$ is large. Whether it is large is a parameter value, which is why the result that $F$ is no
better than $G$ is accommodated rather than implied.

\paragraph{[P] Numbers.} The 57 fields of \texttt{scripts/28} (\texttt{data/interim/field\_vs\_generic.csv})
have mean $r_{FG}=0.786$, mean $c_F=0.401$ and mean $c_G=0.450$.
\begin{itemize}
\item Over fields, $\mathrm{sr}_F$ averages 0.055 (median 0.074), against a threshold that averages 0.141 (median
0.144).
\item $\mathrm{sr}_F$ exceeds the threshold in 15 of 57 fields. By the identity, these are the same fields in which
$c_F>c_G$.
\item At the means, the threshold is 0.156 and $\mathrm{sr}_F$ is 0.077.
\item The raw advantage $c_F-c_G=-0.049$ splits exactly into a dilution term, mean $(r_{FG}c_G-c_G)=-0.083$, and a
department channel, mean $(c_F-r_{FG}c_G)=+0.034$. The channel is of the same order as the mean partial
correlation of $F$ with pay net of $G$ and selectivity ($+0.059$ \ci{+0.003}{+0.114}, \texttt{scripts/58}).
\end{itemize}

\paragraph{[P] Reliability-corrected form.} On true scores,
$\mathrm{Corr}(F^*,G^*)=r_{FG}/\sqrt{\mathrm{rel}_F\,\mathrm{rel}_G}$, which is 0.81--0.88 across the reliability
bracket. With $\mathrm{rel}_Y=0.78$, $\rho^*_G\approx0.51$. Write $F^*=\cos\vartheta^*G^*+\sin\vartheta^*H$, with $H$
the department-specific part of academic prestige. Then $F^*$ out-predicts $G^*$ if and only if
$\rho_{HT}>\rho^*_G\tan(\vartheta^*/2)\approx0.13$--$0.17$. The reliability-corrected advantage is $-0.005$
\ci{-0.043}{+0.044} under the public-edge reliability and $-0.042$ \ci{-0.077}{-0.001} under the extrapolated one
(\texttt{scripts/56}). At the point estimates, therefore, the true-score semi-partial correlation of the
department-specific part with pay is about 0.13--0.17, close to the threshold; the upper end of the interval
allows more. That is not negligible. The department-specific part is diluted by the name, not shown to be absent.
Under the lower-bound reliability, the disattenuated mean partial correlation is $+0.159$ \ci{+0.052}{+0.266} and
the within-institution slope $+0.0297$, about 35\% of the slope without institution fixed effects
(\texttt{scripts/58}).

\paragraph{Agreement of orderings, not of values [S].} Let both valuations be linear indices of a latent
attribute vector $z_i$ (research standing, graduates' ability, resources, location): $s=\alpha'z$ and $m=\nu'z$.
Then
\[
r(F,Y)=\sqrt{\mathrm{rel}_F\,\mathrm{rel}_Y}\;\cos_\Sigma(\alpha,\nu),\qquad
\cos_\Sigma(\alpha,\nu)=\frac{\alpha'\Sigma\nu}{\sqrt{\alpha'\Sigma\alpha\;\nu'\Sigma\nu}} ,
\]
where $\Sigma$ is the covariance of attributes among the institutions that offer the field. If one factor dominates
$\Sigma$ and both $\alpha$ and $\nu$ give it non-zero weight, then $\cos_\Sigma\to\pm1$: audiences that value
different attributes still agree on orderings. With $\mathrm{rel}_Y=0.78$ and $\mathrm{rel}_F\in[0.81,0.96]$, the
mean coupling of $+0.42$ ($r\approx0.44$ under A0) corresponds to $\cos_\Sigma\approx0.50$--$0.55$. Coupling also
depends on $\Sigma$, that is, on which institutions offer the field.

\paragraph{One-factor corollary [B].} Suppose $z=\lambda w+\delta$ with $\mathrm{Var}(w)$ dominant, and suppose the
academy's and selectivity's loadings on $w$ ($\rho_{Aw}$, $\rho_{Sw}$, both non-zero) are common across fields,
while the field's pay loading $\rho_{Yw,f}$ is proportional to $\beta_f$. Then
\[
r_f(F,Y)\approx\sqrt{\mathrm{rel}_F\mathrm{rel}_Y}\,\rho_{Aw}\,\rho_{Yw,f},\qquad
r_f(\text{SAT},Y)\approx\sqrt{\mathrm{rel}_Y}\,\rho_{Sw}\,\rho_{Yw,f}.
\]
Both are proportional to the field's reward for the institution factor, and their true-score correlation across
fields is one. The equal-$\beta$ version of the wage equation cannot produce this; it needs $\beta_f$. The
near-collinearity of the coupling map with field-specific selectivity pricing (true-score $r=+0.94$ in the US,
Spearman $+0.92$ across UK subjects) is what one institution factor implies under this added assumption. Two
consequences follow.
\begin{itemize}
\item The field map is a map of how each field's labour market pays for institution-wide status. It would look
the same with any index that loads on $w$ (faculty salaries, spending, endowment), so for the paper's headline
object (the heterogeneity of $\rho_f$) the question of what $F$ measures does not arise. Test T5 checks the
invariance; its selectivity version is already known.
\item The construct reading of $F$ matters for the claims below the name: the within-institution remainder, $F$
against $G$, and T2.
\end{itemize}

\paragraph{Selectivity.} Average SAT and the admission rate are outcomes of two-sided sorting between applicants
and admission offices. We read them as noisy measures of the selection part of $x_i$, the ability of an
institution's intake; in \citet{macleod2015reputation} and \citet{macleod2017big} a school's reputation is the
expected ability of its graduates. \citet{avery2013revealed} build a revealed-preference ranking of colleges from
applicants' matriculation choices precisely because admission and yield rates are manipulable and poor measures of
applicants' preferences, so their ranking is an analogy for how choices reveal an ordering, not a source for
reading SAT or the admission rate as one. Partialling selectivity out removes the part of the market's valuation
that is measured intake ability. What remains bundles value added, peers, networks, growth and unmeasured
ability. For the question ``does the academy's standing carry information that the market prices beyond
intake?'', selectivity is a confound; for the question ``where do the two orderings agree?'', it locates the
agreement. The numbers (\texttt{scripts/55}):
\begin{itemize}
\item SAT and admission rate alone take mean coupling from $+0.43$ to $+0.14$ (47 fields).
\item The broad set, which adds Pell share, control and the state earnings level, gives $+0.14$ (46 fields), and
$+0.13$ with state fixed effects (28 fields). Controls dated to the graduates' own entry years give the same
($+0.134$ against $+0.129$ on common programs; \texttt{scripts/68}).
\item Pell share, control and the state earnings level alone take it to $+0.31$ (47 fields). Raw $\rho_f$
therefore contains location and composition, which lie outside both constructs.
\end{itemize}
Partial couplings are upper bounds on agreement beyond measured intake only if $F$ and pay load with the same
sign on true selectivity; noisy proxies then leave part of that direction in the residual.

% -------------------------------------------------------------------------------------
\subsection{Proposition 3: below the name}\label{si:theory:umbrella}

\paragraph{Statement and proof [S].} Under A-orth, the department term of (\ref{eq:T-decomp}) is
$\Lambda_{ft}\kappa_f\sigma_r\sigma_y$, with $\Lambda_{ft}=\omega_f+(1-\omega_f)\lambda_t\ge\lambda_t$.
\begin{enumerate}
\item[(i)] At entry with $b=0$, $\lambda_0=0$ and the term is $\omega_f\kappa_f\sigma_r\sigma_y$. It is non-zero if
and only if $\kappa_f\neq0$ and $\omega_f>0$. Only here does attribution matter.
\item[(ii)] For $t>0$, or for $b>0$, $\Lambda_{ft}\ge\lambda_t>0$. Each graduate's own record prices $y_{if}$
whether or not employers track departments, and the term is non-zero if and only if $\kappa_f\neq0$.
\end{enumerate}
The school term behaves the same way with $\kappa^S_d$ and $\omega^S_d$. $\square$

\paragraph{Calibration.} With $b=0$ and $k=3$, the order of magnitude that \citet{lange2007speed} estimates (half
of employers' initial expectation errors gone within about three years), $\lambda_1=0.25$, $\lambda_4=0.57$ and
$\lambda_{10}=0.77$. For college graduates, \citet{arcidiacono2010beyond} find that ability is largely revealed at
hiring, that is, $b$ is large; with $b=0.5$ and $k=3$, $\lambda_4=0.79$. Under either, $\Lambda_{f4}\ge0.57$, and
a department term of about 1\% of 4-year pay per SD of department prestige implies that $\kappa_f\sigma_y$ is
small: department standing carries little information about graduates' market productivity (content). The
information reading (small $\omega_f$) can account for a small 4-year remainder only if learning about individuals
is slow over the horizons observed, $b\approx0$ and $k\gtrsim20$ ($\lambda_4\le0.17$, $\lambda_{10}\le0.33$), which
conflicts with both studies. Information versus content is not identified by the results in hand, and no
registered test separates them.

\paragraph{Dynamics [S].} $\partial\Lambda_{ft}/\partial t=(1-\omega_f)(1-b)\,k/(t+k)^2\ge0$. Learning moves the
department loading away from zero if and only if $\kappa_f\neq0$, $\omega_f<1$ and $b<1$; it raises the loading
only if, in addition, $\kappa_f>0$ (with $\kappa_f<0$, as chemistry's negative remainder would require, it pushes
the loading further below zero). A flat department loading implies $\kappa_f\approx0$, $\omega_f\approx1$, or
$\lambda_t$ nearly constant over the horizons observed ($b\approx1$, or negligible learning).

\paragraph{Umbrella ordering [S].} Suppose learning is the only source of change ($g\equiv0$, no change in who is
observed) and $\omega_i\ge\omega_f$, and take the name prior $\mathbb E[x\mid P_i]$ to be uncorrelated with $a$ (a
prior that covaries positively with $a$ only slows the rise of the institution loading, which strengthens the
inequality). Write $b_I(t)=\beta_f[\omega_i+(1-\omega_i)\lambda_t]\chi\sigma_a\sigma_x$ and
$b_D(t)=\Lambda_{ft}\kappa_f\sigma_r\sigma_y$. Per unit of content, the department loading then rises at least as
fast as the institution loading:
\[
\frac{\dot b_D}{\kappa_f\sigma_r\sigma_y}=(1-\omega_f)\lambda'_t\ \ge\ (1-\omega_i)\lambda'_t=\frac{\dot b_I}{\beta_f\chi\sigma_a\sigma_x} .
\]

\paragraph{[P] The entry pattern (in tension).} On the PSEO fixed cohorts the field-prestige coefficient is
$b_F=+0.189$ \ci{+0.060}{+0.309}, $+0.176$ and $+0.177$ at y1, y5 and y10, while the academia-wide coefficient is
$b_G=-0.013$ \ci{-0.121}{+0.092}, $+0.166$ and $+0.328$ (\texttt{scripts/52}). In the model, with $G$ nearly free of
error ($\mathrm{rel}_G\approx0.99$) and $F$ loading on $a$ and $r$, the joint regression gives $b_F\approx$ the
department loading scaled by the reliability of $F$'s department-specific part, and $b_G\approx$ the name-plus-growth
loading minus $\alpha_f b_F$. Learning alone produces the pattern for no parameter values:
\begin{itemize}
\item $\kappa_f\approx0$ implies $b_F\approx0$, against $b_F>0$;
\item $\kappa_f>0$ with $\omega_f<1$ and $\lambda_t$ rising implies a rising $b_F$, against a flat one;
\item $\omega_f=1$ implies $\omega_i=1$ (Lemma~B, with comparable prior variances), so under learning alone the
name loading ($b_G+\alpha_fb_F$) is flat, against a rising one.
\end{itemize}
The rise of the name loading needs growth, migration, composition or a scale effect in any case
(Proposition~4b). Four readings remain for the flat, positive $b_F$, and the data do not choose among them:
\begin{enumerate}
\item employers know departments from entry ($\omega_f\approx1$, $\kappa_f>0$), against the umbrella of Lemma~B;
\item $\lambda_t$ is nearly constant over y1--y10 ($\kappa_f>0$), either because ability is revealed at hiring
\citep{arcidiacono2010beyond} or because learning is negligible;
\item $F$ tracks a trait that is not the department term, for example a field-specific reward for institution
traits that $F$ measures better than $G$, or program-level selection;
\item y1 is an artefact: M4 fails most at y1, where graduate study and coverage that falls with prestige depress
the medians of high-status institutions, which pushes $b_G$ towards zero.
\end{enumerate}
Net of selectivity the pattern weakens ($b_F=+0.162$ at y1 and $+0.092$ at y10 in \texttt{scripts/66}, sample A),
and in the pooled form, or with $F$ corrected for reliability, $G$'s lead in the rise is not distinguishable from
zero ($+0.0219$ \ci{-0.0067}{+0.0504}; \texttt{scripts/67}). Because M4 fails at y1, the entry-horizon reading of
pay as employer reputation (the Entry corollary) cannot be checked on these data.

\paragraph{[P] Where the remainder appears.} Computer science has the one individually clear within-institution
remainder ($+0.050$ \ci{+0.012}{+0.096}). Engineering and business lean positive; the other 28 fields average zero
or below. This grouping was attached after \texttt{scripts/58}'s results were seen, and it carries little
theoretical weight once Proposition~3(ii) holds: after entry, a remainder requires only $\kappa_f\neq0$, so it says
where department standing carries information about graduates' productivity, not where employers track
departments. Special education ($+0.139$, 15 programs) fits neither story well, and chemistry's negative remainder
($-0.083$) requires $\kappa_f<0$. Whether the computer-science remainder sits in programs that restrict entry is
undetermined: restricted minus open $+0.037$ \ci{-0.051}{+0.146}, with a minimum detectable difference of 0.141
(\texttt{scripts/74}).

% -------------------------------------------------------------------------------------
\subsection{Proposition 4: experience}\label{si:theory:experience}

\paragraph{(a) Decay requires imprecise name beliefs [S].} By Lemma~A, if $\omega_i=1$ no institution-level
loading changes through learning about individuals. With $\omega_i<1$, (\ref{eq:T-wage}) contains
$\beta_f[\omega_i+(1-\omega_i)\lambda_t]x_i+\beta_f(1-\omega_i)(1-\lambda_t)\,\mathbb E[x\mid P_i]$. Regress $m$ on
$x$ (measured by selectivity) and a status measure $G$. The coefficient on $G$ given $x$ is
$\beta_f(1-\omega_i)(1-\lambda_t)\phi_G$, where $\phi_G$ is the slope of the entry prior $\mathbb E[x\mid P]$ on $G$
holding $x$ fixed: the status component of the prior that $G$ tracks beyond selectivity. The coefficient on $x$
given $G$ is $\beta_f[\omega_i+(1-\omega_i)\lambda_t]$. Hence:
\begin{itemize}
\item the $G$ loading net of selectivity decays as $\lambda_t$ rises if and only if $\omega_i<1$ and $\phi_G\neq0$,
whether $\phi_G$ comes from rational shrinkage towards a status-based prior or from a halo; this is the signature
of \citet{altonji2001employer};
\item the selectivity loading net of $G$ rises with $\lambda_t$ if and only if $\omega_i<1$.
\end{itemize}
The results fit as follows.
\begin{itemize}
\item \emph{No decay.} None of 9 decay tests in \texttt{scripts/66} has an interval below 0; in partial form the
status loading net of selectivity rises by $+0.014$ \ci{+0.003}{+0.024} per year, and in the UK it is flat
($+0.008$). This is consistent with $\omega_i\approx1$ or with $\phi_G\approx0$. It says nothing about learning
about individuals (Lemma~A). It bears on the model only insofar as $G$ tracks the public signal $P$ on which
employers' prior rests; the model does not assume that employers see $G$.
\item \emph{Rising selectivity loading.} Net of $G$, the selectivity loading rises in the UK by $+0.024$
\ci{+0.011}{+0.037} per year (regression) and $+0.024$ \ci{+0.016}{+0.032} (partial form), and in the US by $+0.013$
\ci{+0.000}{+0.025} (partial form, at the edge of 0) and $+0.013$ \ci{-0.006}{+0.032} (regression). With
$\omega_i\approx1$, learning cannot produce this.
\item \emph{Together.} A flat status loading and a rising selectivity loading are what learning about names
($\omega_i<1$) with $\phi_G\approx0$ produces. With precisely known names they need selectivity-correlated growth.
Either way they cannot be read as evidence for umbrella precision. One caveat applies to all of these: the loadings are standardised, the
within-cell $R^2$ rises (0.29 to 0.41 in the UK, 0.32 to 0.48 in the US), and a uniform rise of every loading with
$R^2$ gives the same UK pattern; in scale-free form neither loading is singled out (\texttt{scripts/66}).
\end{itemize}

\paragraph{(b) Sources of a rise [B].} In log points, and with $\hat x_i$ as above,
\[
\frac{d}{dt}\mathrm{Cov}(a,m_t)=\beta_f(1-\omega_i)\lambda'_t\big[\chi\sigma_a\sigma_x-\mathrm{Cov}(a,\mathbb E[x\mid P])\big]
+\frac{d}{dt}\mathrm{Cov}(a,g_t)+\frac{d}{dt}\mathrm{Cov}(a,\ell_t)+(\text{composition}),
\]
where $\ell_t$ is the price level and wage premium of where graduates work at $t$. A rising institution-level
covariance therefore requires at least one of four sources.
\begin{itemize}
\item \emph{Status-correlated growth.} With $\Gamma(t)=\mathrm{Cov}(a,g_t)$ rising, for example through campuses
recruited by firms with internal promotion ladders \citep{gibbons1999theory,rivera2015pedigree,weinstein2018employer}.
\citet{macleod2017big} report that college reputation correlates with graduates' earnings growth. Growth need not
be linear.
\item \emph{Name learning.} It raises the covariance only if $\mathrm{Cov}(a,\mathbb E[x\mid P])<\chi\sigma_a\sigma_x$,
that is, if employers' prior under-weights status; with a halo it lowers it. The effect scales with
$1-\omega_i$, so it is largest where employers have seen few graduates.
\item \emph{Migration.} Graduates of high-status institutions may move toward high-wage labour markets over their
careers, which raises pay through price levels as well as productivity.
\item \emph{Composition.} The observed population can change, for example as graduate students enter the
earnings data (a failure of M4). A coverage control removes 19\% of the US rise (\texttt{scripts/52}); the
institution's PhD rate absorbs 0.62 \ci{0.41}{0.83} of it, but that rate also tracks prestige (\texttt{scripts/67}).
\end{itemize}
On standardised scales, which are the scales of every rise reported, a fifth source enters: the loading rises when
the variance of pay unrelated to status shrinks, for example as entry-year local-market components wash out. The
only log-point version (\texttt{scripts/66}, exploratory, net of selectivity) is $+0.0022$ \ci{-0.0005}{+0.0049}
log points per SD per year. Net of SAT and admission rate dated to the graduates' own entry, the rise keeps 0.39 of
its size ($+0.0112$ \ci{+0.0026}{+0.0198} per year; \texttt{scripts/68}).

\paragraph{[P] Shape, conditional on linear growth.} The fixed-cohort coupling rises by $+0.034$
\ci{+0.016}{+0.052} per year from y1 to y5 and by $+0.028$ \ci{+0.017}{+0.037} per year from y5 to y10, a ratio of
about 1.2 (\texttt{scripts/52}). If growth is linear, pure learning with $b=0$ implies the benchmarks below. The
first two columns follow from $k$ alone; the last assumes that learning produces the whole y1--y10 rise of about
0.28.

\medskip
\begin{center}
\small
\begin{tabular}{@{}lccc@{}}
\toprule
learning speed & segment ratio & $\lambda_5-[\tfrac59\lambda_1+\tfrac49\lambda_{10}]$ &
curvature if learning gives the whole rise\\
\midrule
$k=3$ (Lange's order) & 3.2 & 0.144 & about $+0.08$\\
$k=10$ & 1.8 & 0.061 & about $+0.04$\\
$k=30$ & 1.3 & 0.014 & about $+0.02$\\
\bottomrule
\end{tabular}
\end{center}
\medskip

\noindent The curvature is $\rho(5)-[\tfrac59\rho(1)+\tfrac49\rho(10)]$. Observed within-window curvatures are
$+0.036$, $-0.002$, $+0.025$ and $+0.011$ for the four graduation windows. The age curvature of the APC
meta-regression is $+0.024$ \ci{-0.013}{+0.062}, identified only under a piecewise-linear period restriction
(\texttt{scripts/65}). Fast learning implies a concave rise, and the near-linear rise argues against fast learning
as its main source. The converse does not hold: experience-earnings profiles are concave, so status-correlated
growth can be concave too, and concavity would not identify learning. Shape therefore cannot separate learning from
growth; only where the rise occurs can (test T3). The evidence is also weak. If learning at $k=3$ produced the whole
rise, the segment slopes would be about $+0.050$ and $+0.016$ per year, the first inside the observed y1--y5
interval and the second just outside the y5--y10 interval. The calculation applies $\lambda_t$ to correlations,
whose denominators change with $t$, rather than to covariances, so it is approximate, and within-cohort changes
include calendar time.

\paragraph{[P] Access at entry, pay later.} In PSEO Flows, on one cohort and without intervals, placement into
market-priced sectors couples with family prestige at $+0.28$ at y1 and $+0.38$ at y10, while wage coupling goes
from $+0.21$ to $+0.51$ (\texttt{scripts/61}, provisional). Placement coupling rises too, but less than wage
coupling. The pattern fits access followed by growth. It does not exclude learning, because entry placement can
rest on employers' priors, and it does not show employer rationing, because graduates' own choices can produce the
same placement.

\paragraph{(c) Entry conditions.} The model has no entry-condition term. A competitive spot market prices each
graduate at expected productivity whatever the state of the labour market at entry. The Great Recession and COVID
windows (\texttt{scripts/65}) are outside the model.

% -------------------------------------------------------------------------------------
\subsection{Proposition 5: pay schedules}\label{si:theory:schedules}

\paragraph{Derivation [S] (to first order, mean of log pay).} Let a share $\pi_p$ of program $p$'s graduates be
paid on a schedule, with mean level $\bar\tau_p$ where they work, and the rest at market pay
$m^{\rm mkt}_p=\mu^{\rm mkt}_f+\beta_fx_i+\dots$. Treating the program's median as its mean log pay (M4; the median
of a mixture is not linear in $\pi$), $m_p=(1-\pi_p)m^{\rm mkt}_p+\pi_p\bar\tau_p$. With
$\pi_p=\bar\pi_f+\tilde\pi_p$, to first order
\[
\mathrm{Cov}(F,m)=(1-\bar\pi_f)\,\mathrm{Cov}(F,m^{\rm mkt})+\bar\pi_f\,\mathrm{Cov}(F,\bar\tau)
-(\overline{m^{\rm mkt}}-\overline{\tau})\,\mathrm{Cov}(F,\tilde\pi) .
\]
The last term is sorting: higher-status programs send fewer graduates into scheduled jobs when market pay exceeds
schedule pay. $\square$

\paragraph{Grade progression.} Under a national schedule, pay is $\tau(\text{grade}_t)$ and
$\bar\tau_p(t)=\sum_g\Pr(\text{grade}_t=g\mid p)\,\tau_g$. At entry every graduate starts on the entry grade, so
$\mathrm{Cov}(F,\bar\tau)=0$. Over careers,
$\mathrm{Cov}(F,\bar\tau_t)=\sum_g\tau_g\,\mathrm{Cov}\{F,\Pr(\text{grade}_t=g\mid p)\}$, which grows if promotion
across grades (NHS bands, teachers' leadership ranges) depends on status. Grade attainment is a quantity that
employers allocate. Exit from scheduled jobs into market jobs, if it depends on status, moves $\tilde\pi$ over
time and has the same effect.

\paragraph{Consequences.}
\begin{enumerate}
\item[(i)] Under a national schedule, coupling at entry is $(1-\bar\pi_f)\mathrm{Cov}_{\rm mkt}$ plus sorting, and
only sorting as $\bar\pi_f\to1$. Over careers, grade progression and exit can make it rise.
\item[(ii)] Under local schedules, coupling runs through where graduates work (location, and rationing by employers
with high schedules). Net of the schedule level at the destination it is zero at entry.
\item[(iii)] Where M3 fails the wage stops revealing the market's valuation. If employers ration jobs and
promotions, their valuation of status acts on quantities: who is hired where and who is promoted. Graduates'
choices can produce the same quantities, so reading them as valuation needs rationing. Test T6 registers the
quantity prediction.
\item[(iv)] The mechanism does not involve dispersion: schedules can be widely dispersed by location and seniority.
Consistently, the licensed share is unrelated to cross-institution pay dispersion at the point estimate ($+0.02$
\ci{-0.30}{+0.32}, 45 fields; \texttt{scripts/35c}), although that interval is too wide to test (iv), and the
within-occupation dispersion law is null (\texttt{scripts/53}).
\item[(v)] The field's reward for institution traits is scaled by the same $(1-\bar\pi_f)$: $\beta_f$ becomes
$(1-\bar\pi_f)\beta_f$. A schedule gradient is therefore collinear with weaker selectivity pricing and cannot be
separated from it in these data.
\end{enumerate}

\paragraph{Results.} UK medicine and nursing couple weakly and are flat over careers (slopes $-0.053$ and $-0.018$
per year; \texttt{scripts/62}), as (i) implies at the point where schedule shares are high, but the schedule share
is coded from outside and coincides with field type. UK teaching rises from $+0.012$ at YAG1 to $+0.211$ at YAG5
($+0.050$ per year). Grade progression and exit accommodate this, but they were added to the model after the rise
was seen. Two results sit badly with Proposition~5:
\begin{itemize}
\item US special education couples at $+0.55$ despite district salary schedules; its partial correlation given
geography is $+0.59$ to $+0.66$, on 19 institutions.
\item UK allied health couples at $+0.570$ on NHS Agenda for Change bands.
\end{itemize}
Proposition~5 accommodates these only through an effective $\pi$ below the nominal one (employment outside the
schedule, allowances, exit), which is ad hoc without sector-level pay data. At the cell level
(\texttt{scripts/61}), the SAT gradient is flatter where more graduates enter setting sectors, but the prestige and
$G$ gradients are not robustly flatter, and the attribution to setting or market is not identified.
Proposition~5's weak-coupling clause cannot fail on wage data by construction; its content is the quantity
prediction of T6.

% -------------------------------------------------------------------------------------
\subsection{Mapping of results to propositions (organised after the results were seen)}\label{si:theory:map}

Fit classes:
\begin{itemize}
\item \textbf{S}: implied by the structure for every admissible parameter value. No result is in this class.
\item \textbf{B}: implied under an added assumption named in the row.
\item \textbf{A}: accommodated by parameter values, most of them inferred from these same results, or by a post hoc
classification.
\item \textbf{N}: not implied; accommodated only by an element added to the model.
\item \textbf{C}: contradicted, or in tension.
\item \textbf{O}: outside the model, or a check on the indicator rather than on the model.
\item \textbf{P}: a post hoc calibration, conditional on a stated assumption.
\end{itemize}

{\small
\begin{longtable}{@{}p{0.47\linewidth}p{0.16\linewidth}p{0.31\linewidth}@{}}
\caption{Every empirical result mapped to a proposition. Numbers are those of the main text and SI; sources in
parentheses.}\label{tab:T-map}\\
\toprule
Result & Proposition & Fit\\
\midrule
\endfirsthead
\toprule
Result & Proposition & Fit\\
\midrule
\endhead
\bottomrule
\endfoot
1. Coupling is positive and heterogeneous: mean $+0.419$ \ci{+0.355}{+0.482} (52 fields; scripts/58); Cochran
$Q=183.2$ on 56 df, $I^2=0.69$ (57 fields; scripts/57); the field ordering replicates on PSEO ($+0.68$; $+0.74$)
(scripts/32, 59) & 2; one-factor corollary & A: positivity follows from $\chi>0$, a primitive; level and spread
come from free field parameters.\\
2. Selectivity absorbs most of the level: SAT and admission rate alone $+0.43\to+0.14$ (47 fields); broad set
$+0.14$ (46); Pell, control and state level alone $+0.31$; era-matched controls give the same (scripts/55, 68).
Shapley $R^2$: institution earnings 0.201, selectivity 0.173, $G$ 0.086, $F$ 0.076 (scripts/58) & 1b; 2 & A: $x$
includes selection, so any absorbed share fits; the part removed by Pell, control and state level lies outside both
constructs.\\
3. The cross-field ordering is collinear with selectivity pricing: true-score $r=+0.94$ (US); Spearman $+0.92$
across UK subjects (scripts/55, 62) & 2, one-factor corollary & B: under one institution factor with a
field-specific $\beta_f$.\\
4. Field prestige is no better than $G$: $c_F-c_G=-0.049$; reliability-corrected $-0.005$ \ci{-0.043}{+0.044} to
$-0.042$ \ci{-0.077}{-0.001} (scripts/28, 56) & 2 (identity, threshold) & A: the identity sets the threshold (0.16
against 0.08 at the means); a $\kappa_f\Lambda_{ft}$ below it is a parameter value.\\
5. Within-institution remainder $+0.0084$ \ci{+0.0000}{+0.0164}, about 10\% of the slope without fixed effects;
mean partial $+0.059$ \ci{+0.003}{+0.114} net of $G$ and selectivity; under the lower-bound reliability $+0.0297$
and $+0.159$ \ci{+0.052}{+0.266} (scripts/58) & 2(iii); 3 & A: after entry it requires only $\kappa_f\neq0$; small
$\kappa_f\sigma_y$, or $b\approx0$ with slow learning and small $\omega_f$; not identified. The within design is
robust to $\mathrm{Cov}(r,x)$ only as far as the traits span $x$.\\
6. The remainder is individually clear only in computer science ($+0.050$ \ci{+0.012}{+0.096}; the one field that
survives a multiple-testing correction); engineering and business lean positive; the remaining 28 fields are zero
or below on average. Uncorrected intervals also exclude 0 for special education ($+0.139$, 15 programs) and
chemistry ($-0.083$) (scripts/58) & 3 & A, post hoc: after entry the remainder marks content ($\kappa_f\neq0$), not
tracking; chemistry needs $\kappa_f<0$; special education fits neither story.\\
7. Computer-science entry policy: restricted minus open $+0.037$ \ci{-0.051}{+0.146}, minimum detectable difference
0.141 (scripts/74) & 3 & O: undetermined (underpowered).\\
8. UK: 66--90\% of coupling survives standardisation on prior attainment. What survives tracks intake selectivity:
within band, $G$ adds $+0.033$ \ci{-0.027}{+0.098} and selectivity $+0.252$ beyond $G$ (scripts/62, 72) & 1b; 2 &
A: consistent in direction (a school-mean premium net of own band); the share is not predicted.\\
9. Coupling rises with years since graduation on fixed cohorts: US $+0.031$/yr \ci{+0.021}{+0.038} (career-time
part $[+0.014,+0.032]$), UK $+0.026$/yr. A coverage control removes 19\%; the institution PhD rate absorbs 0.62
\ci{0.41}{0.83}; net of era-matched SAT and admission rate the slope keeps 0.39 of its size (scripts/52, 59, 62,
67, 68) & 4(b) & N: needs growth, migration, composition or a scale effect.\\
10. The rise loads on $G$, not $F$: $dG=+0.038$ \ci{+0.023}{+0.051}, $dF=-0.001$ \ci{-0.017}{+0.013} (scripts/52,
59); in pooled form, or with $F$ corrected for reliability, $G$'s lead is not distinguishable from 0 (scripts/67)
& 3 dynamics; 4(b) & A: the $G$ rise is the added growth of row 9; the flat $F$ loading is row 11.\\
11. At y1 $F$ carries the association and $G$ does not: $b_F=+0.189$ \ci{+0.060}{+0.309}, $b_G=-0.013$
\ci{-0.121}{+0.092}; $b_F$ stays flat ($+0.176$, $+0.177$) (scripts/52) & 3 dynamics & C: learning alone fits no
parameter values; left are known departments (against Lemma~B), a nearly constant $\lambda_t$, a non-departmental
trait tracked by $F$, or a y1 artefact (M4).\\
12. No decay of the status loading net of selectivity: 0 of 9 tests; partial form $+0.014$ \ci{+0.003}{+0.024}/yr;
UK $+0.008$/yr (scripts/66) & 4(a) & A: $\omega_i\approx1$ or $\phi_G\approx0$; uninformative about learning about
individuals.\\
13. The selectivity loading net of $G$ rises: UK $+0.024$ \ci{+0.011}{+0.037} (regression), $+0.024$
\ci{+0.016}{+0.032} (partial); US partial $+0.013$ \ci{+0.000}{+0.025} (edge), regression $+0.013$
\ci{-0.006}{+0.032}; a uniform rise of standardised loadings with $R^2$ gives the UK pattern (scripts/66) & 4(a) &
A: name learning with $\omega_i<1$ and $\phi_G\approx0$, selectivity-correlated growth, or a scale effect; with
row 12, it does not favour umbrella precision.\\
14. Shape of the rise: segments $+0.034$ and $+0.028$/yr; age curvature $+0.024$ \ci{-0.013}{+0.062} under a
restriction (scripts/52, 65) & 4(b) & P: if growth is linear, fast learning is not the main source; concavity
would not separate learning from concave growth.\\
15. Program-median reliability is flat across horizons (0.756/0.736/0.740) (scripts/73) & M4 & O: a check on the
indicator. The model does not imply it: learning about individuals widens within-program dispersion.\\
16. Placement into market sectors couples with prestige as pay does ($+0.40$ vs $+0.41$); higher-status programs
send fewer graduates into setting-priced sectors ($\rho=-0.19$ \ci{-0.27}{-0.11}). Placement coupling rises from
$+0.28$ (y1) to $+0.38$ (y10) while pay coupling rises from $+0.21$ to $+0.51$ (one cohort). Within institutions,
one SD of department prestige goes with $+0.51$ pp of market-sector share against $+3.33$ pp across institutions
(scripts/61, provisional) & 3; 4(b); 5(iii) & A: reads as status acting on quantities only if employers ration
sector access; graduates' sorting predicts the same (T8).\\
17. Pay-scale subjects (UK nursing, medicine, teaching) couple at $+0.120$ against $+0.523$; careers are flat in
medicine ($-0.053$/yr) and nursing ($-0.018$/yr); the US nursing premium is the state's RN wage (scripts/50, 62) &
5(i)--(ii) & A: given $\pi_f$, which is coded from outside and coincides with field type.\\
18. UK teaching rises from $+0.012$ (YAG1) to $+0.211$ (YAG5), $+0.050$/yr (scripts/62) & 5, grade progression & N:
grade progression and exit were added after the rise was seen (T6).\\
19. The US strict-licensure gradient ($-0.53$ to $-0.48$) runs only between disciplines and cannot be separated from
field type or from selectivity pricing (scripts/64) & 5(v) & A: the collinearity is implied, so the contrast cannot
identify the mechanism.\\
20. US special education ($+0.55$) and UK allied health ($+0.570$) couple despite pay schedules (scripts/62, 64) &
5 & C: accommodated only through an ad hoc effective $\pi$.\\
21. Cell-level Flows: the SAT gradient is flatter with the setting share; the prestige and $G$ gradients are not
robustly flatter; attribution is not identified (scripts/61) & 5, cell level & C (tension).\\
22. Location at the institution's own state: the field-specific state earnings component moves mean coupling from
$+0.419$ to $+0.415$ and leaves the within-institution slope at $+0.0084$ (scripts/69) & 1b (location) & A: covers
the origin state only; the destination is untested (T7).\\
23. Great Recession entry window: $D=-0.064$ \ci{-0.114}{-0.014}, opposite to the pre-specified sign (placebo rank
$p=0.20$); the state-unemployment interaction is fragile (scripts/65) & none & O: no entry-condition term.\\
24. Cross-window drift of $+0.015$/yr at fixed horizon, with window-to-window swings (scripts/65) & none & O:
reputations are static in the model.\\
25. Elite tail: $G$ adds little beyond selectivity and parental income, and no more for the top 1\% than for the
median (scripts/60, provisional) & none & O: the model is median-based. Mild tension with an access reading if
ladders concentrate at the top.\\
26. The Pell-share loading falls from $+0.243$ at y1 to $+0.011$ at y10; PSEO coverage falls with prestige at y1
(scripts/52, 66) & M4 & O: composition, a failure of the indicator, not valuation.\\
27. Coupling rises with field size; after the earnings-noise correction the log-$n$ slope falls from $+0.077$ to
$+0.027$ (scripts/57, 73) & 2 (reliability; range of $\Sigma$) & A: the range term is untested.\\
28. Disciplinary kinship does not organise coupling (Moran's $I=+0.011$, $p=0.95$, 14 fields; 30\% power at a kin
correlation of 0.30) (scripts/51) & none & O.\\
\end{longtable}
}

% -------------------------------------------------------------------------------------
\subsection{Registered tests (not yet run)}\label{si:theory:tests}

\paragraph{Status.} The tests were specified on 25 September 2026, before any statistic named in them was
computed, in two binding texts: the original specification B0 (T1--T3) and an amendment A1 written the same day
after an adversarial review of B0 (T0, T4--T8, and additions to T1--T3). None has been run. The tests share data
with results already seen. T0 and T1 use the Wapman edges, analysed with split-half SpringRank in
\texttt{scripts/56}. T2 uses $F$, $G$ and Scorecard pay from \texttt{scripts/28} and \texttt{56}, and its pay
contrast is already known. T3 uses the PSEO panel of \texttt{scripts/66} and was chosen knowing that the $G$
loading rises. T4 uses the UK band panels of \texttt{scripts/62}, where the within-band rise of raw $G$ coupling
($+0.017$ per year) is known. T5 has a known selectivity analogue, and T6--T8 use PSEO Flows, whose national
sector shares were analysed in \texttt{scripts/61}. The dates are self-attested. The binding texts are kept
verbatim in \theoryprereg, and the script checks their SHA-256 digests before running:
\begin{center}\footnotesize
B0: \texttt{bfaa08a283f6ebb5aa96face4d3e8768e2911572066750588d97210083243441}\\
A1: \texttt{5e1b9d7ec132a29a9a722c1bebf86b49f94cd3b46cb95dff06f80122d2a498c1}
\end{center}
A digest detects later edits; it does not establish when a text was written. This subsection summarises the two
texts; where they differ from it, the binding texts govern. Every verdict will be reported, whatever it is.

Before B0 was written we made only outcome-blind checks, none of which computes a statistic named in the tests:
\begin{itemize}
\item every public edge has weight at least 2;
\item the share of persons without recorded gender. B0 states 9.4\%, which is the share pooled over the academia,
domain and field networks, which count the same persons several times. In the field-level edges that T1 uses it
is 6.2\%, and 5.3\% after dropping self-hires (recomputed before A1 was written; A1 records the erratum);
\item no script uses the gender columns;
\item the size of the ORCID sample. Over all 681 shards, 968 (economics) to 10{,}112 (chemistry) persons have a US
bachelor's episode and a later US doctorate at a different institution in the same keyword field, and each field
has 100--701 bachelor's-institution cells with at least three persons before matching to Wapman institutions. No
destination rank or earnings figure was computed.
\item the presence of the needed columns.
\end{itemize}

\paragraph{Standing of the two texts.} B0 is not edited; its verdicts are computed as registered, with
clarifications that resolve ambiguities in it (for example, that its T3 interval is the normal two-way interval of
\texttt{scripts/59}, and that a postdoctoral role is not a doctorate). A1 adds statistics with their own verdicts
and a claim map that fixes, before any data are seen, which verdict governs each claim in the paper. Both verdicts
are always reported.

\paragraph{T0: production audit of $F$ (Proposition 1a, H0).} Per field, the Spearman correlations of $F$ with
log PhD production (IPEDS research doctorates in the field's CIP codes), with the export ratio
$(d^{\rm out}+1)/(d^{\rm in}+1)$, with the placement rate (out-degree per doctorate) and with Wapman's
ProductionRank; and of $F$ with a minimum-violation ordering found by local search on the public edges. Then
coupling with ranks from a degree-corrected Bradley--Terry fit (offsets $\log(d^{\rm out}_ud^{\rm in}_v/
d^{\rm out}_vd^{\rm in}_u)$) and with $F$ residualised on production and the export ratio, as shares of coupling
with $F$. \emph{Rule}: production-robust if both retained shares have lower bounds of at least 0.75;
production-dominated if either upper bound is below 0.50, in which case $F$ is renamed ``placement position''.
\emph{Rivals.} Placement power predicts a low retained share and valuation a high one, so T0 discriminates between
them; it does not separate valuation from other network positions.

\paragraph{T1: single-index audit (Proposition 1a).} B0: for each institution, $P_u$ (mean field percentile of the
employers of its PhDs) and $H_u$ (mean percentile of its faculty's PhD origins); $R$ is the ratio of cross-role to
within-role split-half Spearman correlations, supported if its lower bound is at least 0.80. B0 also compares
SpringRank fitted separately to men's and women's hires against a permutation null. A1 adds a cross-fitted $R'$
(ranks fitted on one half of persons, $P$ and $H$ computed on the other), benchmarked against $R_0$, the value that
SpringRank's own generative model produces on simulated networks with the fitted ranks, $\beta$, and the public
censoring: one index if the lower bound of $R'-R_0$ is at least $-0.10$. A1 treats the gender comparison as an audit
of H3 (a uniform gender penalty leaves both orderings unchanged) and adds an equivalence rule: invariant only if the
upper 95\% bound of the mean Spearman shortfall is at most 0.10, with the minimum detectable effect reported
first; otherwise uninformative. \emph{Rivals.} Network position and placement power predict the same agreement of
roles, so T1 does not discriminate between them and valuation, and a strongly hierarchical network almost always
passes T1(a): it is a low-risk audit. A failed T1 withdraws only the reading of $F$ as a
belief about the quality of a department's PhDs.

\paragraph{T2: the academy's choice among the same graduates (Propositions 1--3).} From the ORCID mobility edges
\citep{li2026orcid}, persons with a US bachelor's episode and a later US research doctorate at a different
institution in the same field (computer science, mathematics, physics, psychology, economics, chemistry, biology).
$D_{if}$ is the mean field percentile of the doctoral destinations of program $(i,f)$'s graduates. B0's statistics
are $\Delta^D=\rho(F,D)-\rho(G,D)$, the partial $\rho(F,D\mid G)$, and $\Delta^Y=\rho(F,Y)-\rho(G,Y)$ on the cells
with pay. $\Delta^Y$ is already known for all seven fields (computer science $-0.045$, mathematics $+0.004$,
physics $-0.052$, psychology $-0.035$, economics $-0.005$, chemistry $-0.073$, biology $-0.091$), so only the $D$
side is new, and B0's primary $\Delta^D$ is dominated by dilution: by the identity it is negative unless the
department channel exceeds $c^D_G(1-r_{FG})$. A1 therefore adds T2$'$: the mean partial correlations
$\rho(F,D\mid G)$ and $\rho(F,D^G\mid G)$, where $D^G$ scores destinations on the academia-wide scale (removing the
advantage $F$ has from sharing $D$'s scale), and their contrast with $\rho(F,Y\mid G)$ on the same cells. Supported
if all three lower bounds exceed 0; contradicted if both $D$ upper bounds are below 0; inconclusive otherwise,
including when both conditions hold. A1 also reports the reliability of $D$, a reliability-corrected $\Delta^D$, the
share of same-institution doctorates by $F$, a control for doctorates in the same state, and word-anchored field
keywords, and it excludes postdoctoral and predoctoral roles that B0's regular expression admits.
\emph{Auxiliary assumption.} The ordering that governs faculty hiring also governs doctoral admission and
applicants' choice of doctoral program; the model itself contains no doctoral admission. \emph{Rivals} that
predict the same sign: advisor-network and subfield routing, geography, content (strong departments prepare
researchers better), self-selection and the shared scale of $F$ and $D$. T2 is therefore a consistency check
with one new prediction, not a double dissociation, and it does not separate information from content.

\paragraph{T3: where the career rise occurs (Proposition 4b).} B0: in each field $\times$ cohort $\times$ horizon
cell of the PSEO panel, the standardised rank of pay is regressed on $G$, visibility $V^\perp$ (log undergraduate
enrolment residualised on $G$) and $G\times V^\perp$; $\theta_{GV}$ is the per-year slope of the interaction.
Learning where employers rarely see an institution's graduates predicts $\theta_{GV}<0$; access at nationally
recruited institutions predicts $\theta_{GV}>0$. A1 adds T3$'$ for three reasons. B0's ``null'' requires a minimum
detectable effect that comparable slopes in \texttt{scripts/66} make nearly unreachable. With a status-leaning
(halo) prior, learning gives $\theta_{GV}>0$. And low enrolment also marks small private institutions that recruit
nationally. T3$'$ adds public/private control and its interaction with $G$; reads $\theta'$ against the sign of the
y1 interaction $g_1$, since learning shrinks the entry gap whatever its sign; reports the simple slopes of the $G$
gradient at low and high visibility; computes the minimum detectable effect by permuting pay ranks within cells
before any estimate; and requires a learning signature to hold in public and in private institutions. The decision
table is: $g_1>0$ and $\theta'<0$, learning; $g_1>0$ and $\theta'>0$, access or growth at high visibility; $g_1<0$
and $\theta'>0$, halo correction or access; $g_1<0$ and $\theta'<0$, growth concentrated at low visibility. A
secondary exposure measure is the share of graduates working out of state (PSEO Flows); learning and national
recruiting predict the same sign on it, so it does not discriminate on its own.

\paragraph{T4: UK within-band decay (employer reputation in the learning sense).} Within prior-attainment bands of
the UK panels (YAG 1, 3, 5), the per-year slope of the partial Spearman correlation of band median pay with
institution selectivity given $G$. Employer learning with graduates' attainment unseen by employers predicts
decay; value added, growth, or employers who see attainment predict none. Decay if the interval lies below 0, rise
if above, null if it includes 0 with a minimum detectable effect of at most 0.02 per year.

\paragraph{T5: invariance of the field map.} For faculty salary, instructional spending per student and endowment
per student, the true-score correlation across fields between each index's coupling map and $F$'s, from split
halves of institutions (fields with at least 60 institutions). The one-factor corollary predicts values near one:
supported if all three lower bounds are at least 0.70; contradicted if two upper bounds are below 0.70.

\paragraph{T6: quantities where pay is scheduled (Proposition 5).} In education and registered nursing, the
Spearman correlation of $G$ with the scheduled-pay level of the census divisions where an institution's graduates
work (PSEO Flows division shares at y5 times ACS median pay of teachers or registered nurses aged 23--35 by
division). Supported if the lower bound exceeds 0; contradicted if the upper bound is below $+0.05$, which would
withdraw the quantity clause of Proposition~5.

\paragraph{T7: destination location.} On the PSEO fixed cohorts, the share of coupling at y5 that survives
partialling out a destination wage index (PSEO Flows division shares times ACS division wages of young bachelor's
holders), and the share of the career rise it removes, alone and with the out-of-state share. Location-robust if the
retained share has a lower bound of at least 0.75; location-dominated if its upper bound is below 0.50.

\paragraph{T8: between- and within-sector coupling.} On the \texttt{scripts/61} cells, coupling with and without the
sector-wage-mix index (Flows sector shares times ACS national sector wages of young bachelor's holders), reported as
a between-sector share with its interval and no verdict.

\paragraph{Claim map.} A1 fixes which claim each verdict governs. T0 governs whether $F$ is named academic prestige
or placement position; T1$'$ governs the single-index reading; T2$'$ governs the consistency claim about the
academy's own choice; T3$'$ governs the source of the rise; T4 governs whether the within-band premium is called
reputation in the learning sense; T5, T6 and T7 govern the map's invariance, the quantity clause of Proposition~5
and whether construct claims about pay remain after destination location is netted out. No verdict licenses a
claim about the mechanism below the name: none separates employers' information from the content of department
standing.

\paragraph{Candidate tests not registered.} Some implications were not registered because their outcome is largely
known or their design is contaminated:
\begin{itemize}
\item a school-level umbrella test with a leave-own-field sibling mean (the engineering and business remainders are
known, and Wapman's domains are not organisational schools: chemical and biomedical engineering sit in its Natural
Sciences domain);
\item a school-level umbrella test with the Wapman domain rank (it shares $F$'s edges and sampling error);
\item a status measure that leaves out the field's own domain (nearly collinear with $G$, whose
selectivity-controlled slope is known);
\item a mechanism test that would separate $\omega_f$ from $\kappa_f$ through fields whose $F$ loading rises with
experience (the per-field slopes of the $F$ loading were computed in \texttt{scripts/52}).
\end{itemize}
Others are feasible and uncontaminated but lower in priority:
\begin{itemize}
\item range restriction (the spread of $G$ among the institutions offering a field as a moderator of coupling);
\item growth in within-program dispersion as a measure of the pay-setting regime that uses neither prestige nor
licensing codes;
\item exit mediation within nursing and education, beyond T6's secondary exit share;
\item convergent validity of $F$ against a direct reputational survey \citep{nrc1995research}.
\end{itemize}

% -------------------------------------------------------------------------------------
\subsection{What the model cannot do}\label{si:theory:limits}

\begin{enumerate}
\item \emph{It is post hoc.} The parameters $\chi$, $\beta_f$, $\kappa_f$, $\omega_f$, $\Gamma$ and $\pi_f$ are free,
and several of them ($\kappa_f$ small for most fields, $\omega_i\approx1$ or $\phi_G\approx0$) are inferred from the
results they organise. No result is implied for every parameter value. Only the registered tests test the model.
\item \emph{It predicts no magnitudes.} It says nothing about the size of $+0.43\to+0.14$, the 66--90\% UK
survival, or a remainder of 1\% of pay per SD.
\item \emph{It does not identify the mechanism below the name.} After entry, a department term requires only
content; information and content are separable only under slow learning about individuals, which the evidence on
employer learning does not support.
\item \emph{Growth is a reduced-form slot.} Networks, skills, promotion ladders, migration and sorting into steeper
careers all enter $\Gamma$ or $\ell_t$. Aggregate data cannot tell them apart, nor rule out very slow learning.
\item \emph{Selection and value added are bundled.} Both enter $x$ and $y$. The model justifies the valuation
reading of coupling but licenses no causal statement about what institutions or departments add.
\item \emph{It uses a single index per audience and a linear-Gaussian form.} It collapses subfield prestige and
sector-specific employer beliefs, cannot represent premia concentrated at the top, and treats medians as means of
log pay.
\item \emph{Reputations are static.} Calendar and cohort drift and entry-condition effects are outside the model;
accommodating them would need dynamic collective reputations \citep{tirole1996theory}.
\item \emph{Scheduled fields.} By design, the wage indicator fails there. How employers of nurses, doctors or
teachers value status can be learned only from quantities (T6).
\item \emph{The UK.} The UK $G$ rests on ORCID edges whose external validity is 0.74--0.92, and the UK has no
department axis, so UK construct claims are weaker than US ones.
\item \emph{The construct defence rests on assumptions.} H0--H4, A-orth and M1--M4 can be made plausible but not
verified pair by pair. T0 and T1 are the audits that could show the reading of $F$ wrong; T7 and T8 bound what the
reading of pay includes.
\end{enumerate}
